\chapter{Sistema proposto}
\label{cap:rad-sistema-proposto}

\section{Panoramica}

Tech4All è un'applicazione web che organizza contenuti didattici in un
catalogo consultabile liberamente e ne verifica l'apprendimento tramite quiz
corretti dal sistema. Il percorso tipico dell'utente è quello di
\cref{fig:rad-ad-apprendimento}: dalla ricerca di un contenuto alla sua
lettura, dalla verifica al riconoscimento del traguardo raggiunto.

Due decisioni caratterizzano il sistema e ricorrono in tutto il documento.

\begin{description}[leftmargin=1.2cm,style=nextline]
  \item[Consultazione libera] Catalogo, tutorial, feedback e quiz sono
    visibili senza autenticazione. La registrazione è richiesta solo per
    \emph{contribuire} (consegnare un quiz, lasciare un feedback) o per
    accedere a dati personali. Questa scelta discende direttamente da
    \id{OB\_1}: chi ha poca dimestichezza con il digitale è spesso scoraggiato
    proprio dalla richiesta di registrarsi.
  \item[Verifica non aggirabile] La correzione di un quiz è un'operazione del
    sistema, non del dispositivo dell'utente. Il quiz consegnato al browser
    non contiene l'indicazione della risposta corretta, che viene rivelata
    solo insieme all'esito. È la condizione perché \id{OB\_2} e \id{OB\_3}
    abbiano senso: un progresso che si può falsificare non misura nulla.
\end{description}

\diagramma{0.78}{ad-apprendimento}{Percorso di apprendimento nel sistema
proposto.}{fig:rad-ad-apprendimento}

\section{Attori}
\label{sec:rad-attori}

Il sistema riconosce tre attori, in relazione di generalizzazione fra loro
(\cref{tab:rad-attori}). Non esistono attori esterni non umani: il sistema non
si integra con servizi di terze parti.

\begin{table}[H]
\centering
\begin{tabular}{@{}L{3.0cm} L{11.2cm}@{}}
\toprule
\textbf{Attore} & \textbf{Descrizione} \\
\midrule
Visitatore &
Chiunque raggiunga la piattaforma senza essere autenticato. Può consultare
il catalogo, leggere i tutorial e i feedback altrui, vedere le domande di un
quiz. Non può contribuire né accedere a dati personali. \\

Utente &
Visitatore che si è registrato e autenticato. Oltre a quanto precede, svolge
i quiz, lascia feedback, consulta i propri progressi e amministra il proprio
account. Generalizza \emph{Visitatore}. \\

Amministratore &
Utente con privilegi di gestione. Pubblica e aggiorna i contenuti, definisce
i quiz e gli obiettivi, modera i feedback e amministra gli account.
Generalizza \emph{Utente}: conserva tutte le capacità di un utente ordinario.
\\
\bottomrule
\end{tabular}
\caption{Attori del sistema.}
\label{tab:rad-attori}
\end{table}

\begin{nota}
Il ruolo di amministratore non è ottenibile dall'interno del sistema: non
esiste alcuna funzionalità di promozione. Il primo amministratore è creato
in fase di inizializzazione della base di dati. È una scelta deliberata:
un'operazione di questa portata non deve poter essere raggiunta attraverso
l'interfaccia applicativa.
\end{nota}

\section{Requisiti funzionali}
\label{sec:rad-rf}

I requisiti funzionali sono raggruppati per area, con il prefisso che ne
identifica il sottosistema di appartenenza:
\id{AU} per l'accesso e gli account,
\id{CT} per il catalogo dei tutorial,
\id{QZ} per i quiz,
\id{FB} per i feedback,
\id{OB} per gli obiettivi.
La formulazione completa di ciascun requisito è riportata
nell'appendice~\ref{cap:rad-app-requisiti}; qui se ne descrive la logica complessiva.

\subsection{Accesso e gestione dell'account}

Un visitatore può registrarsi fornendo nome, cognome, email e password
(\id{RF\_AU\_1}) e autenticarsi con le sole credenziali (\id{RF\_AU\_2}). La
disponibilità di un'email può essere verificata prima dell'invio del modulo
(\id{RF\_AU\_5}), così che l'utente non scopra il conflitto solo a
registrazione tentata.

L'utente autenticato consulta il proprio profilo (\id{RF\_AU\_4}), ne
modifica i dati anagrafici (\id{RF\_AU\_6}), cambia la propria password
esibendo quella attuale (\id{RF\_AU\_7}), chiude la sessione
(\id{RF\_AU\_3}) e può eliminare definitivamente il proprio account
(\id{RF\_AU\_8}). L'amministratore consulta l'elenco degli utenti registrati
(\id{RF\_AU\_9}) e può eliminarne l'account (\id{RF\_AU\_10}).

\subsection{Catalogo dei tutorial}

Il catalogo è consultabile per intero (\id{RF\_CT\_1}), ristretto a una
categoria (\id{RF\_CT\_2}) o interrogato per parola chiave su titolo, testo e
categoria (\id{RF\_CT\_3}); il singolo tutorial è consultabile per
identificativo (\id{RF\_CT\_4}). L'elenco delle categorie ammesse è esposto
dal sistema (\id{RF\_CT\_5}) anziché replicato nell'interfaccia: è la stessa
sorgente a definire i valori accettati e quelli proposti.

L'amministratore pubblica un tutorial (\id{RF\_CT\_6}), lo aggiorna
(\id{RF\_CT\_7}) o lo ritira (\id{RF\_CT\_8}). Le immagini inserite nel corpo
del contenuto sono caricate (\id{RF\_CT\_9}) e rimosse (\id{RF\_CT\_10}) come
operazioni distinte, perché il loro ciclo di vita è quello della redazione,
non quello del tutorial.

\subsection{Quiz}

Il quiz di un tutorial è consultabile da chiunque, ma privo delle soluzioni
(\id{RF\_QZ\_1}). L'utente autenticato consegna le proprie risposte e riceve
esito, numero di risposte esatte, soluzioni ed eventuali badge sbloccati
(\id{RF\_QZ\_2}). L'amministratore definisce le domande di un quiz
(\id{RF\_QZ\_3}), le ridefinisce integralmente (\id{RF\_QZ\_4}) o rimuove il
quiz (\id{RF\_QZ\_5}).

\subsection{Feedback}

I feedback di un tutorial sono pubblici (\id{RF\_FB\_1}); l'utente consulta i
propri (\id{RF\_FB\_2}), ne inserisce al più uno per tutorial
(\id{RF\_FB\_3}), lo modifica (\id{RF\_FB\_4}) o lo elimina (\id{RF\_FB\_5}).
L'amministratore può eliminare il feedback di un altro utente a fini di
moderazione (\id{RF\_FB\_6}).

\subsection{Obiettivi e badge}

Gli obiettivi raggiungibili sono pubblici (\id{RF\_OB\_1}), così che l'utente
sappia quali traguardi lo attendono; i badge conseguiti sono consultabili
nella propria area (\id{RF\_OB\_2}). L'assegnazione è automatica e avviene
come conseguenza del superamento di un quiz (\id{RF\_OB\_3}): non esiste
alcuna funzionalità per assegnare un badge a mano. L'amministratore definisce
(\id{RF\_OB\_4}), aggiorna (\id{RF\_OB\_5}) ed elimina (\id{RF\_OB\_6}) gli
obiettivi.

\section{Requisiti non funzionali}
\label{sec:rad-rnf}

I requisiti non funzionali seguono il modello FURPS+, integrato con le
categorie di ISO/IEC 25010~\cite{iso25010}. Ogni requisito è accompagnato da
un criterio di accettazione osservabile: requisiti che non si possono
verificare non sono stati inclusi.

\subsection{Funzionalità (sicurezza)}

\begin{longtable}{@{}l L{5.2cm} L{6.6cm}@{}}
\toprule
\textbf{Id} & \textbf{Requisito} & \textbf{Criterio di accettazione} \\
\midrule
\endfirsthead
\toprule
\textbf{Id} & \textbf{Requisito} & \textbf{Criterio di accettazione} \\
\midrule
\endhead
\bottomrule
\endlastfoot

\id{RNF\_F\_1} &
Le password sono conservate esclusivamente come hash. &
Nessuna risposta del sistema contiene la password o il suo hash; il valore
memorizzato è un hash bcrypt~\cite{provos1999}. \\

\id{RNF\_F\_2} &
Nessuna operazione di modifica è eseguibile senza una sessione valida. &
Ogni richiesta di modifica priva di sessione riceve risposta \id{401}. \\

\id{RNF\_F\_3} &
Le operazioni di gestione sono riservate agli amministratori. &
Le stesse richieste, da un utente ordinario autenticato, ricevono \id{403}. \\

\id{RNF\_F\_4} &
L'autore di un'operazione è determinato dalla sessione, non dalla richiesta. &
Un identificativo utente inviato nel corpo di una richiesta viene ignorato. \\

\id{RNF\_F\_5} &
Il contenuto dei tutorial non può veicolare codice eseguibile. &
Marcatori \id{script}, attributi di evento e URL con schema \id{javascript:}
non sono presenti nel contenuto memorizzato~\cite{owasp2021}. \\

\id{RNF\_F\_6} &
Le soluzioni di un quiz non sono conoscibili prima della consegna. &
La rappresentazione di un quiz non contiene alcun campo che distingua la
risposta corretta. \\

\id{RNF\_F\_7} &
I nomi di file forniti dall'esterno non permettono di uscire dalla cartella
consentita. &
Un nome contenente separatori di percorso o risalite viene rifiutato con
\id{400}. \\

\end{longtable}

\subsection{Usabilità}

\begin{longtable}{@{}l L{5.2cm} L{6.6cm}@{}}
\toprule
\textbf{Id} & \textbf{Requisito} & \textbf{Criterio di accettazione} \\
\midrule
\endfirsthead
\toprule
\textbf{Id} & \textbf{Requisito} & \textbf{Criterio di accettazione} \\
\midrule
\endhead
\bottomrule
\endlastfoot

\id{RNF\_U\_1} &
L'interfaccia usa un linguaggio non tecnico e in lingua italiana. &
Nessun messaggio rivolto all'utente contiene termini informatici non
definiti nel glossario. \\

\id{RNF\_U\_2} &
Ogni errore è comunicato indicando che cosa correggere. &
Le risposte di validazione elencano i campi non conformi con un messaggio in
linguaggio naturale. \\

\id{RNF\_U\_3} &
Le operazioni irreversibili richiedono una conferma esplicita. &
Eliminazione di account, tutorial, quiz e feedback presentano una richiesta
di conferma prima di essere inviate. \\

\id{RNF\_U\_4} &
La consultazione dei contenuti non richiede registrazione. &
Catalogo, tutorial, feedback, quiz e obiettivi rispondono \id{200} a un
chiamante non autenticato. \\

\end{longtable}

\subsection{Affidabilità}

\begin{longtable}{@{}l L{5.2cm} L{6.6cm}@{}}
\toprule
\textbf{Id} & \textbf{Requisito} & \textbf{Criterio di accettazione} \\
\midrule
\endfirsthead
\toprule
\textbf{Id} & \textbf{Requisito} & \textbf{Criterio di accettazione} \\
\midrule
\endhead
\bottomrule
\endlastfoot

\id{RNF\_A\_1} &
Le operazioni che coinvolgono più entità sono atomiche. &
L'interruzione della creazione di un quiz non lascia quiz privi di domande. \\

\id{RNF\_A\_2} &
Un errore imprevisto non rivela dettagli interni del sistema. &
La risposta \id{500} non contiene messaggi del DBMS, percorsi di file o
tracce di esecuzione. \\

\id{RNF\_A\_3} &
Il comportamento del sistema è verificato automaticamente. &
La copertura dei test non scende sotto l'85\% delle istruzioni e l'80\% delle
diramazioni; la soglia è imposta dalla suite, non solo dichiarata. \\

\id{RNF\_A\_4} &
Il conteggio dei progressi resta corretto anche in caso di tentativi
ripetuti. &
Ripetere un quiz già superato non altera il numero di quiz superati. \\

\end{longtable}

\subsection{Prestazioni}

\begin{longtable}{@{}l L{5.2cm} L{6.6cm}@{}}
\toprule
\textbf{Id} & \textbf{Requisito} & \textbf{Criterio di accettazione} \\
\midrule
\endfirsthead
\toprule
\textbf{Id} & \textbf{Requisito} & \textbf{Criterio di accettazione} \\
\midrule
\endhead
\bottomrule
\endlastfoot

\id{RNF\_P\_1} &
Il caricamento di un quiz non degrada con il numero di domande. &
L'aggregato quiz-domande-risposte è letto con una sola interrogazione. \\

\id{RNF\_P\_2} &
Ricerca e filtro sono eseguiti dalla base di dati. &
Il sistema non trasferisce l'intero catalogo per poi ridurlo altrove. \\

\end{longtable}

\subsection{Supportabilità}

\begin{longtable}{@{}l L{5.2cm} L{6.6cm}@{}}
\toprule
\textbf{Id} & \textbf{Requisito} & \textbf{Criterio di accettazione} \\
\midrule
\endfirsthead
\toprule
\textbf{Id} & \textbf{Requisito} & \textbf{Criterio di accettazione} \\
\midrule
\endhead
\bottomrule
\endlastfoot

\id{RNF\_S\_1} &
I sottosistemi condividono la stessa organizzazione interna. &
Ogni sottosistema è composto da rotte, servizio, DAO ed entità con le
medesime responsabilità. \\

\id{RNF\_S\_2} &
I vincoli di dominio sono definiti una sola volta. &
Lunghezze, formati e soglie hanno un'unica definizione, condivisa fra
validazione delle richieste e regole di dominio. \\

\id{RNF\_S\_3} &
Ogni modifica al codice è verificata automaticamente. &
Analisi statica, controllo dei tipi, test e compilazione sono eseguiti a
ogni integrazione. \\

\end{longtable}

\section{Vincoli}
\label{sec:rad-vincoli}

I vincoli, o pseudo-requisiti, non descrivono ciò che il sistema fa ma entro
quali limiti deve essere realizzato.

\begin{table}[H]
\centering
\begin{tabular}{@{}l L{12.2cm}@{}}
\toprule
\textbf{Id} & \textbf{Vincolo} \\
\midrule
\id{V\_1} & Il sistema è un'applicazione web, raggiungibile con un browser e
            senza installazione sul dispositivo dell'utente. \\
\id{V\_2} & La comunicazione fra interfaccia e sistema avviene tramite
            un'interfaccia HTTP con rappresentazioni JSON. \\
\id{V\_3} & La persistenza è affidata a un DBMS relazionale. \\
\id{V\_4} & Il trattamento dei dati personali è conforme al GDPR: l'utente
            può consultare, rettificare ed eliminare i propri dati. \\
\id{V\_5} & Il sistema è realizzato con tecnologie open source e non dipende
            da servizi a pagamento per il proprio funzionamento. \\
\bottomrule
\end{tabular}
\caption{Vincoli di progetto.}
\label{tab:rad-vincoli}
\end{table}

\section{Funzionalità escluse dal perimetro}
\label{sec:rad-fuori-perimetro}

Tre funzionalità sono state deliberatamente escluse. Sono elencate qui perché
un requisito assente è un'informazione tanto quanto uno presente, e perché la
loro assenza è una scelta e non una dimenticanza.

\begin{description}[leftmargin=1.2cm,style=nextline]
  \item[Assistente conversazionale] Il sistema può incorporare un widget di
    dialogo fornito da terzi, configurabile e disattivabile. Non implementa
    logica conversazionale propria: realizzarla richiederebbe competenze e
    infrastrutture estranee agli obiettivi del progetto, e un assistente
    inaffidabile su un pubblico inesperto sarebbe dannoso più che utile.
  \item[Recupero della password] Richiede l'invio di posta elettronica
    transazionale, quindi un servizio esterno e un dominio verificato. In sua
    assenza, un utente che perde le credenziali deve registrarsi nuovamente.
  \item[Supporto multilingua] Interfaccia e contenuti sono in italiano. Il
    sistema non è predisposto per la traduzione dei tutorial, che sono testi
    redatti e non stringhe di interfaccia.
\end{description}
